\documentclass[11pt]{article}

% Plantilla común para material de ayudantías de Campos y Ondas.
% Documento A4, una columna, fuente Computer Modern de 11 puntos
% y márgenes de 2.5 cm por lado, en el mismo estilo visual usado
% para las guías de Física Matemática II.

\usepackage[utf8]{inputenc}
\usepackage[T1]{fontenc}
\usepackage[spanish]{babel}
\usepackage{amsmath,amssymb,mathtools,bm,mathrsfs,graphicx,array,enumitem,caption,float,microtype}
\usepackage[dvipsnames]{xcolor}
\usepackage[a4paper,margin=2.5cm]{geometry}
\usepackage{tikz}
\usetikzlibrary{arrows.meta,calc,patterns,babel,decorations.pathreplacing,decorations.pathmorphing,positioning}
\usepackage[
    colorlinks=true,
    linkcolor=red,
    citecolor=green,
    urlcolor=red
]{hyperref}
\spanishdecimal{.}

\setlength{\parindent}{0pt}
\setlength{\parskip}{0.45em}
\setlength{\abovedisplayskip}{6pt}
\setlength{\belowdisplayskip}{6pt}
\setlength{\abovedisplayshortskip}{4pt}
\setlength{\belowdisplayshortskip}{4pt}
\setlist[itemize]{topsep=4pt,itemsep=2pt,leftmargin=1.8em}
\setlist[enumerate]{topsep=4pt,itemsep=2pt,leftmargin=1.8em}
\captionsetup{font=small,labelfont=bf}

% Colores comunes usados en figuras y recuadros.
\definecolor{headinggray}{RGB}{55,55,55}
\definecolor{rulegray}{RGB}{150,150,150}
\definecolor{bluevec}{RGB}{29,78,216}
\definecolor{orangevec}{RGB}{180,83,9}
\definecolor{redvec}{RGB}{185,28,28}
\definecolor{greenvec}{RGB}{21,128,61}
\definecolor{fieldblue}{RGB}{30,90,180}
\definecolor{currentorange}{RGB}{195,95,20}
\definecolor{inducedgreen}{RGB}{20,120,80}
\definecolor{wiregray}{RGB}{55,55,55}
\definecolor{waveblue}{RGB}{29,78,216}
\definecolor{waveorange}{RGB}{180,83,9}
\definecolor{wavegreen}{RGB}{21,128,61}
\definecolor{wavered}{RGB}{185,28,28}

% Encabezado homogéneo.
\newcommand{\encabezado}[3]{%
{\Large\bfseries #1\par}
\vspace{2pt}
{\normalsize #2\hfill #3\par}
\vspace{6pt}\hrule\vspace{8pt}}

\newcommand{\problema}[2]{%
\vspace{0.55em}
\noindent\textbf{#1.}\enspace{\small #2}\par
\vspace{0.3em}}

\newcommand{\inciso}[1]{\vspace{0.35em}\noindent\textbf{(#1)}\enspace}

% Comandos auxiliares.
\newcommand{\dd}{\,\mathrm{d}}
\newcommand{\ii}{\mathrm{i}}
\newcommand{\e}{\mathrm{e}}
\newcommand{\sen}{\operatorname{sen}}
\newcommand{\grad}{\nabla}
\newcommand{\laplaciano}{\nabla^2}
\newcommand{\R}{\mathbb{R}}

% Comandos usados en las guías de Campos y Ondas.
\newcommand{\soluciontitulo}{%
  \par\medskip
  \noindent\textbf{Solución.}\par
  \smallskip
}

\newcommand{\resultadofinal}[1]{%
  \par\medskip
  \noindent\textbf{Resultado final:} #1\par
}


\newcommand{\paso}[1]{%
  \par\medskip
  \noindent\textbf{#1}\par
  \vspace{0.1em}%
}

\newcommand{\Bentra}[2]{%
    \draw[fieldblue,line width=0.55pt] (#1,#2) circle (0.11);
    \draw[fieldblue,line width=0.55pt] ($(#1,#2)+(-0.065,-0.065)$)--($(#1,#2)+(0.065,0.065)$);
    \draw[fieldblue,line width=0.55pt] ($(#1,#2)+(-0.065,0.065)$)--($(#1,#2)+(0.065,-0.065)$);
}
\newcommand{\Bsale}[2]{%
    \draw[fieldblue,line width=0.55pt] (#1,#2) circle (0.11);
    \fill[fieldblue] (#1,#2) circle (0.03);
}

\tikzset{
    >=Latex,
    eje/.style={->,thick,headinggray},
    onda/.style={very thick,waveblue},
    ondados/.style={very thick,waveorange},
    ondaaux/.style={thick,wavegreen},
    vec/.style={->,very thick,wavered},
    vecV/.style={->,very thick,bluevec},
    vecB/.style={->,very thick,orangevec},
    vecF/.style={->,very thick,redvec},
    vecE/.style={->,very thick,greenvec},
    panel/.style={draw=rulegray!80,fill=black!2,rounded corners=2pt},
    cargaPos/.style={circle,draw=bluevec,fill=blue!8,thick,minimum size=8mm,inner sep=0pt},
    cargaNeg/.style={circle,draw=greenvec,fill=green!8,thick,minimum size=8mm,inner sep=0pt},
    guia/.style={densely dashed,rulegray},
    etiqueta/.style={font=\footnotesize},
    nota/.style={draw=rulegray,fill=white,rounded corners=1.5pt,align=left,font=\footnotesize,inner xsep=6pt,inner ysep=5pt,line width=0.35pt}
}


\newcommand{\Fsale}[2]{%
    \draw[redvec,very thick] (#1,#2) circle (0.16);
    \fill[redvec] (#1,#2) circle (0.04);
}
\newcommand{\Fentra}[2]{%
    \draw[redvec,very thick] (#1,#2) circle (0.16);
    \draw[redvec,very thick] ($(#1,#2)+(-0.095,-0.095)$)--($(#1,#2)+(0.095,0.095)$);
    \draw[redvec,very thick] ($(#1,#2)+(-0.095,0.095)$)--($(#1,#2)+(0.095,-0.095)$);
}

\begin{document}

\encabezado{Inducción electromagnética y ley de Faraday}{Campos y Ondas 510245}{J.R., F.B.}
\shorthandoff{<>}

\section*{Problema 1.}
Un im\'an se mueve hacia una espira como indica la figura. Encuentre la direcci\'on de la corriente a trav\'es de la resistencia:
\begin{enumerate}[label=(\alph*)]
    \item mientras el im\'an cae hacia la espira;
    \item despu\'es de que el im\'an ha pasado a trav\'es de la espira y se aleja de ella.
\end{enumerate}


Este problema no requiere calcular la fem. La pregunta es s\'olo por el sentido de la corriente inducida. Por eso se debe identificar primero si el flujo externo aumenta o disminuye, y luego aplicar la ley de Lenz.


\begin{figure}[H]
\centering
\begin{tikzpicture}[x=0.82cm,y=0.82cm,>=Latex,font=\small]
    % Figura hecha con coordenadas internas m\'as compactas y separadas.
    % No se usa escalado global de una figura grande.

    % Paneles
    \draw[rulegray!50] (-0.45,-1.35) rectangle (5.05,4.60);
    \draw[rulegray!50] (6.10,-1.35) rectangle (11.60,4.60);

    % ---------------- Panel A ----------------
    \begin{scope}[shift={(2.30,0)}]
        \node[font=\bfseries] at (0,4.20) {(a) El im\'an se acerca};

        % Im\'an: polo sur mirando a la espira
        \draw[wiregray,thick,fill=gray!8] (-0.48,2.95) rectangle (0.48,3.75);
        \node at (0,3.55) {$N$};
        \node at (0,3.12) {$S$};
        \draw[fieldblue,very thick,->] (1.05,3.75)--(1.05,2.95)
            node[midway,right] {mov.};

        % Flujo externo ascendente
        \node[fieldblue,align=center] at (0,2.42) {$\Phi_B^{\rm ext}$ ascendente\\\textbf{aumenta}};

        % Espira
        \draw[wiregray,very thick] (0,0.70) circle (1.05);
        \draw[currentorange,very thick,->] (-0.74,1.45) arc (135:-70:1.05);

        % Campo externo saliendo del plano: anillo de 8 simbolos
        \foreach \x/\y in {0.620/0.700,0.438/1.138,0/1.320,-0.438/1.138,-0.620/0.700,-0.438/0.262,0/0.080,0.438/0.262}{\Bsale{\x}{\y}}

        % Campo inducido entrando: 2 simbolos centrales, separados
        \foreach \x in {-0.21,0.21}{
            \draw[inducedgreen,line width=0.95pt] ($(\x,0.70)+(-0.10,-0.10)$)--($(\x,0.70)+(0.10,0.10)$);
            \draw[inducedgreen,line width=0.95pt] ($(\x,0.70)+(-0.10,0.10)$)--($(\x,0.70)+(0.10,-0.10)$);
        }

        \node[inducedgreen,align=center] at (0,-0.70) {$\vec B_{\rm ind}$ entra};
        \node[currentorange,align=center] at (0,-1.10) {$I_{\rm ind}$ horario};
    \end{scope}

    % ---------------- Panel B ----------------
    \begin{scope}[shift={(8.85,0)}]
        \node[font=\bfseries] at (0,4.20) {(b) El im\'an se aleja};

        \draw[wiregray,thick,fill=gray!8] (-0.48,2.95) rectangle (0.48,3.75);
        \node at (0,3.55) {$N$};
        \node at (0,3.12) {$S$};
        \draw[fieldblue,very thick,->] (1.05,2.95)--(1.05,3.75)
            node[midway,right] {mov.};

        \node[fieldblue,align=center] at (0,2.42) {$\Phi_B^{\rm ext}$ ascendente\\\textbf{disminuye}};

        \draw[wiregray,very thick] (0,0.70) circle (1.05);
        \draw[currentorange,very thick,->] (0.74,1.45) arc (45:250:1.05);

        % Campo externo saliendo del plano: anillo de 8 simbolos
        \foreach \x/\y in {0.620/0.700,0.438/1.138,0/1.320,-0.438/1.138,-0.620/0.700,-0.438/0.262,0/0.080,0.438/0.262}{\Bsale{\x}{\y}}

        % Campo inducido saliendo: 2 simbolos centrales, separados
        \foreach \x in {-0.21,0.21}{
            \draw[inducedgreen,line width=0.75pt] (\x,0.70) circle (0.12);
            \fill[inducedgreen] (\x,0.70) circle (0.040);
        }

        \node[inducedgreen,align=center] at (0,-0.70) {$\vec B_{\rm ind}$ sale};
        \node[currentorange,align=center] at (0,-1.10) {$I_{\rm ind}$ antihorario};
    \end{scope}
\end{tikzpicture}
\caption{Vista superior de la espira. El s\'imbolo $\odot$ representa campo saliendo del plano y $\times$ representa campo entrando al plano. La corriente inducida se escoge para oponerse al \textbf{cambio} de flujo, no al campo externo por s\'i solo.}
\label{fig:p1_lenz}
\end{figure}


\soluciontitulo

\paso{Caso (a): el im\'an cae hacia la espira.}
El polo sur del im\'an se acerca a la espira. En la regi\'on de la espira, el campo magn\'etico externo producido por el im\'an apunta hacia arriba. Como el im\'an se acerca, ese flujo ascendente aumenta:
\begin{equation*}
    \Phi_B^{\rm ext}\ \text{ascendente aumenta}.
\end{equation*}

Por la ley de Lenz, la espira induce un campo que \textbf{se opone a ese aumento}. Luego,
\begin{equation*}
    \vec B_{\rm ind}\ \text{debe apuntar hacia abajo}.
\end{equation*}

Usando la regla de la mano derecha, un campo inducido hacia abajo corresponde a una \textbf{corriente horaria} vista desde arriba de la espira. Por tanto,
\begin{equation*}
    \boxed{I_{\rm ind}\ \text{circula en sentido horario}.}
\end{equation*}

\paso{Caso (b): el im\'an ya atraves\'o la espira y se aleja.}
Despu\'es de atravesar la espira, el flujo externo ascendente disminuye:
\begin{equation*}
    \Phi_B^{\rm ext}\ \text{ascendente disminuye}.
\end{equation*}

La espira intenta sostener ese flujo. Por eso produce un \textbf{campo inducido ascendente}:
\begin{equation*}
    \vec B_{\rm ind}\ \text{debe apuntar hacia arriba}.
\end{equation*}

Un campo inducido hacia arriba corresponde a una \textbf{corriente antihoraria} vista desde arriba. As\'i,
\begin{equation*}
    \boxed{I_{\rm ind}\ \text{circula en sentido antihorario}.}
\end{equation*}


La frase correcta no es ``la espira se opone al im\'an''. Lo correcto es: la espira se opone al cambio del flujo magn\'etico que atraviesa su superficie.


\clearpage

\section*{Problema 2.}
Encuentre la direcci\'on de la corriente en la resistencia de la bobina secundaria:
\begin{enumerate}[label=(\alph*)]
    \item en el instante en que se cierra el interruptor;
    \item despu\'es de que el interruptor ha estado cerrado durante varios minutos;
    \item en el instante en que se abre el interruptor.
\end{enumerate}


La bobina secundaria no est\'a conectada a la bater\'ia. Por tanto, cualquier corriente en la secundaria aparece por inducci\'on mutua. La pregunta central es si el flujo magn\'etico que atraviesa la bobina secundaria est\'a cambiando.


\begin{figure}[H]
\centering
\begin{tikzpicture}[x=0.88cm,y=0.88cm,>=Latex,font=\small]
    % Dos bobinas con objetos reubicados, no s\'olo escalados.

    % N\'ucleos
    \draw[wiregray,thick] (0.25,3.00) rectangle (3.95,3.36);
    \draw[wiregray,thick] (6.35,3.00) rectangle (10.05,3.36);

    % Bobinas
    \draw[wiregray,very thick,decorate,decoration={coil,aspect=0.42,segment length=7.5pt,amplitude=7.2pt}] (0.95,3.36)--(3.25,3.36);
    \draw[wiregray,very thick,decorate,decoration={coil,aspect=0.42,segment length=7.5pt,amplitude=7.2pt}] (7.05,3.36)--(9.35,3.36);

    % Circuito primario
    \draw[wiregray,thick] (0.95,3.02)--(0.95,1.25);
    \draw[wiregray,thick] (3.25,3.02)--(3.25,1.25);
    \draw[wiregray,thick] (0.95,1.25)--(1.50,1.25);
    \draw[wiregray,thick] (2.52,1.25)--(3.25,1.25);

    % Interruptor
    \draw[wiregray,fill=white,thick] (1.50,1.25) circle (0.06);
    \draw[wiregray,fill=white,thick] (2.10,1.25) circle (0.06);
    \draw[wiregray,thick] (1.56,1.31)--(2.00,1.82);
    \node[left] at (1.42,1.50) {$S$};
    \draw[wiregray,thick] (2.10,1.25)--(2.35,1.25);

    % Bater\'ia
    \draw[wiregray,thick] (2.35,0.86)--(2.35,1.64);
    \draw[wiregray,thick] (2.72,0.68)--(2.72,1.82);
    \node at (2.54,0.32) {$\mathcal E$};
    \node[align=center] at (2.10,-0.16) {bobina primaria};

    % Circuito secundario
    \draw[wiregray,thick] (7.05,3.02)--(7.05,1.12);
    \draw[wiregray,thick] (9.35,3.02)--(9.35,1.12);
    \draw[wiregray,thick] (7.05,1.12)--(7.55,1.12);
    \draw[wiregray,thick] (8.85,1.12)--(9.35,1.12);
    \draw[wiregray,thick] (7.55,1.12)--(7.70,1.43)--(7.85,0.82)--(8.00,1.43)--(8.15,0.82)--(8.30,1.43)--(8.45,0.82)--(8.60,1.43)--(8.75,0.82)--(8.85,1.12);
    \node at (8.20,0.32) {$R$};
    \node[align=center] at (8.20,-0.16) {bobina secundaria};

    % Flujo acoplado: centrado y con espacio
    \draw[fieldblue,very thick,->] (4.35,3.18)--(5.95,3.18);
    \node[fieldblue,above] at (5.15,3.28) {flujo acoplado};

    % Corrientes inducidas, puestas bajo la bobina para no chocar con el resistor
    \draw[currentorange,very thick,->] (9.20,2.55)--(7.25,2.55);
    \node[currentorange,right] at (9.28,2.55) {al cerrar};
    \draw[inducedgreen,very thick,->] (7.25,1.90)--(9.20,1.90);
    \node[inducedgreen,right] at (9.28,1.90) {al abrir};
\end{tikzpicture}
\caption{Inducci\'on mutua entre bobinas. La separaci\'on espacial permite distinguir el circuito primario, el flujo acoplado y la corriente inducida en la resistencia secundaria.}
\label{fig:p2_mutua}
\end{figure}


\soluciontitulo

\paso{Caso (a): instante en que se cierra el interruptor.}
Al cerrar el interruptor, la corriente en la bobina primaria \textbf{crece desde cero} hasta su valor estacionario:
\begin{equation*}
    I_p:0\longrightarrow I_0.
\end{equation*}

Al crecer $I_p$, tambi\'en crece el campo magn\'etico producido por la primaria. Por tanto, el flujo magn\'etico que atraviesa la secundaria aumenta:
\begin{equation*}
    \Phi_{B,p}\ \text{aumenta en la bobina secundaria}.
\end{equation*}

La bobina secundaria induce una corriente cuyo campo se opone a ese aumento. Con el sentido de enrollamiento mostrado en el enunciado, esto corresponde a corriente por la resistencia de derecha a izquierda:
\begin{equation*}
    \boxed{I_s\ \text{por la resistencia de derecha a izquierda}.}
\end{equation*}

\paso{Caso (b): interruptor cerrado durante varios minutos.}
Despu\'es de un tiempo largo, la corriente primaria \textbf{ya no cambia}:
\begin{equation*}
    I_p=\text{constante}.
\end{equation*}

Entonces el flujo magn\'etico que atraviesa la secundaria tambi\'en es constante:
\begin{equation*}
    \frac{d\Phi_B}{dt}=0.
\end{equation*}

Por la ley de Faraday,
\begin{equation*}
    \varepsilon_s=-N_s\frac{d\Phi_B}{dt}=0.
\end{equation*}

Luego, no hay corriente inducida:
\begin{equation*}
    \boxed{I_s=0.}
\end{equation*}

\paso{Caso (c): instante en que se abre el interruptor.}
Al abrir el interruptor, la corriente primaria cae a cero:
\begin{equation*}
    I_p:I_0\longrightarrow 0.
\end{equation*}

El flujo magn\'etico producido por la primaria disminuye. La secundaria intenta sostener ese flujo, de modo que el sentido de la corriente inducida se invierte respecto del caso de cierre:
\begin{equation*}
    \boxed{I_s\ \text{por la resistencia de izquierda a derecha}.}
\end{equation*}


Un campo magn\'etico constante no induce corriente. Lo que induce corriente es un flujo magn\'etico variable.


\clearpage

\section*{Problema 3.}
Una espira plana conductora que consta de una sola vuelta, con un \textit{\'area} de secci\'on transversal de $8.0\,\mathrm{cm^2}$, es perpendicular a un campo magn\'etico cuya magnitud aumenta uniformemente de $0.5\,\mathrm{T}$ a $2.5\,\mathrm{T}$ en $1.0\,\mathrm{s}$. Determine la corriente inducida resultante si la espira tiene una resistencia de $2.0\,\Omega$.


Aqu\'i el campo magn\'etico es uniforme y perpendicular a la espira. Por eso el flujo se escribe directamente como $\Phi_B=BA$. No hace falta integrar.


\begin{figure}[H]
\centering
\begin{tikzpicture}[x=0.90cm,y=0.90cm,>=Latex,font=\small]
    \draw[very thick] (0,0) circle (1.05);
    \draw[currentorange,thick,->] (0,1.05) arc (90:405:1.05);
    \foreach \x/\y in {-0.45/-0.45,0/-0.45,0.45/-0.45,-0.45/0,0/0,0.45/0,-0.25/0.45,0.25/0.45}{\Bentra{\x}{\y}}
    \node[align=center] at (0,-1.70) {campo uniforme\\perpendicular};
    \draw[fieldblue,very thick,->] (1.75,-0.82)--(3.75,-0.82) node[midway,below] {$\Delta t=1.0\,\mathrm{s}$};
    \node[align=center] at (2.75,0.18) {$B_i=0.5\,\mathrm{T}$\\$B_f=2.5\,\mathrm{T}$};
    \draw[<->] (-1.05,1.34)--(1.05,1.34) node[midway,above] {$A=8.0\,\mathrm{cm^2}$};
\end{tikzpicture}
\caption{Espira de una vuelta en un campo magn\'etico uniforme cuya magnitud cambia linealmente con el tiempo.}
\label{fig:p3_uniforme}
\end{figure}

\soluciontitulo

Los datos son
\begin{equation*}
    A=8.0\,\mathrm{cm^2},\qquad
    B_i=0.5\,\mathrm{T},\qquad
    B_f=2.5\,\mathrm{T},\qquad
    \Delta t=1.0\,\mathrm{s},\qquad
    R=2.0\,\Omega.
\end{equation*}

Primero se convierte el \textit{\'area} a \textbf{unidades SI}:
\begin{equation*}
    A=8.0\,\mathrm{cm^2}=8.0\times10^{-4}\,\mathrm{m^2}.
\end{equation*}

Como el campo es \textbf{perpendicular} a la espira,
\begin{equation*}
    \Phi_B=BA.
\end{equation*}

La magnitud de la fem inducida es
\begin{align*}
    |\varepsilon|
    &=\left|\frac{\Delta\Phi_B}{\Delta t}\right| \\
    &=A\left|\frac{\Delta B}{\Delta t}\right|.
\end{align*}

La variaci\'on del campo magn\'etico es
\begin{equation*}
    \Delta B=B_f-B_i=2.5\,\mathrm{T}-0.5\,\mathrm{T}=2.0\,\mathrm{T}.
\end{equation*}

Luego,
\begin{align*}
    |\varepsilon|
    &=(8.0\times10^{-4})\frac{2.0}{1.0}\,\mathrm{V} \\
    &=1.6\times10^{-3}\,\mathrm{V}.
\end{align*}

La corriente inducida se obtiene con la \textbf{ley de Ohm}:
\begin{align*}
    I_{\rm ind}
    &=\frac{|\varepsilon|}{R} \\
    &=\frac{1.6\times10^{-3}}{2.0}\,\mathrm{A} \\
    &=8.0\times10^{-4}\,\mathrm{A} \\
    &=0.8\,\mathrm{mA}.
\end{align*}


El signo de Faraday se usa para fijar el sentido de la corriente inducida. Como este problema s\'olo pide la corriente resultante, basta entregar la magnitud.


\clearpage

\section*{Problema 4.}
Una espira rectangular de ancho $w$ y largo $L$ y un alambre largo y recto que transporta una corriente se encuentran sobre una mesa como se muestra en la figura.
\begin{enumerate}[label=(\alph*)]
    \item Determine el flujo magn\'etico a trav\'es de la espira debido a la corriente $I$.
    \item Suponga que la corriente cambia con el tiempo seg\'un $I=a+bt$, donde $a$ y $b$ son constantes. Determine la fem inducida si $b=10\,\mathrm{A/s}$, $h=1.0\,\mathrm{cm}$, $w=10.0\,\mathrm{cm}$ y $L=100\,\mathrm{cm}$. Indique la direcci\'on de la corriente inducida.
\end{enumerate}


Este problema \textbf{no se resuelve con $\Phi_B=BA$} porque el campo del alambre no es uniforme. Cerca del alambre el campo es mayor y lejos del alambre es menor. Por eso se integra el flujo sobre la superficie de la espira.


\begin{figure}[H]
\centering
\begin{tikzpicture}[x=0.72cm,y=0.72cm,>=Latex,font=\small]
    % Figura del problema 4 corregida desde coordenadas internas.
    % r se mide desde el alambre hasta la cara superior de la franja.
    % La cara inferior queda a distancia r+dr.

    % Coordenadas principales
    \coordinate (wireL) at (0.00,6.20);
    \coordinate (wireR) at (15.45,6.20);
    \coordinate (loopTL) at (4.20,3.60);
    \coordinate (loopBR) at (13.10,1.00);
    \coordinate (stripTL) at (4.20,2.60);
    \coordinate (stripBR) at (13.10,2.15);

    % Alambre largo y corriente
    \draw[wiregray,very thick] (wireL)--(wireR);
    \draw[currentorange,very thick,->] (6.30,6.20)--(8.90,6.20)
        node[midway,above=3pt] {$I(t)$};

    % Espira rectangular
    \draw[wiregray,very thick] (loopTL) rectangle (loopBR);

    % Campo magnetico entrando al plano. Mas denso cerca del alambre (arriba).
    \foreach \x in {4.65,5.45,6.25,7.05,7.85,8.65,9.45,10.25,11.05,11.85,12.65}{\Bentra{\x}{3.10}}
    \foreach \x in {4.90,6.30,7.70,9.10,10.50,11.90}{\Bentra{\x}{1.85}}
    \foreach \x in {5.30,6.90,11.30,12.50}{\Bentra{\x}{1.25}}
    \node[wiregray] at (8.80,1.25) {espira};

    % Franja diferencial dA = L dr
    \fill[inducedgreen!10] (stripTL) rectangle (stripBR);
    \draw[inducedgreen!70!black,very thick] (stripTL) rectangle (stripBR);
    \node[inducedgreen!70!black,fill=white,inner sep=1.5pt] at (8.65,2.375) {$dA=L\,dr$};

    % Espesor dr de la franja, indicado en el borde derecho.
    \draw[inducedgreen!70!black,thick,<->] (13.65,2.60)--(13.65,2.15)
        node[midway,right,fill=white,inner sep=1pt] {$dr$};

    % Dimensiones de la espira.
    \draw[<->,thick] (4.20,0.45)--(13.10,0.45) node[midway,below] {$L$};
    \draw[<->,thick] (14.75,3.60)--(14.75,1.00) node[midway,right] {$w$};

    % Lineas guia (lineas de extensi\'on hacia las flechas h y r).
    \draw[dashed,rulegray] (1.30,3.60)--(4.20,3.60);
    \draw[dashed,rulegray] (0.40,2.60)--(4.20,2.60);

    % h: distancia del alambre al borde superior de la espira.
    \draw[<->,thick] (1.30,6.20)--(1.30,3.60) node[midway,left] {$h$};

    % r: distancia del alambre a la cara superior de la franja.
    \draw[<->,thick] (0.40,6.20)--(0.40,2.60) node[midway,left] {$r$};

    % Identificacion de las dos caras de la franja: r y r+dr.
    \node[anchor=east,fill=white,inner sep=1pt] at (4.10,2.60) {cara en $r$};
    \node[anchor=east,fill=white,inner sep=1pt] at (4.10,2.15) {cara en $r+dr$};

    % Notas fisicas, en zona libre.
    \node[align=left] at (6.80,4.85) {variable de integraci\'on:\\[-1mm] $h\le r\le h+w$};
    \node[align=left] at (11.40,4.85) {$B(r)=\dfrac{\mu_0 I}{2\pi r}$};
\end{tikzpicture}
\caption{Construcci\'on del elemento de \textit{\'area} para el flujo. La coordenada $r$ se mide desde el alambre hasta la cara superior de la franja. La cara inferior est\'a a distancia $r+dr$, por lo que el espesor de la franja es $dr$ y su \textit{\'area} es $dA=L\,dr$. La variable de integraci\'on es $r\in[h,h+w]$.}
\label{fig:p4_alambre}
\end{figure}


\soluciontitulo

\paso{Parte (a): flujo magn\'etico.}
El campo magn\'etico de un \textbf{alambre largo y recto} a una distancia $r$ es
\begin{equation*}
    B(r)=\frac{\mu_0 I}{2\pi r}.
\end{equation*}

Se toma una \textbf{franja diferencial horizontal} de largo $L$ y espesor $dr$, paralela al alambre. La coordenada que recorre la espira es $r$; \textbf{no se escribe $dr\in[h,h+w]$}. El diferencial $dr$ s\'olo representa el espesor infinitesimal de la franja. Entonces
\begin{equation*}
    dA=L\,dr.
\end{equation*}

El flujo diferencial a trav\'es de esa franja es
\begin{align*}
    d\Phi_B
    &=B(r)dA \\
    &=\frac{\mu_0 I}{2\pi r}L\,dr.
\end{align*}

Como la espira empieza a distancia $h$ del alambre y termina a distancia $h+w$, se integra usando \textbf{$r$ como variable de integraci\'on} en el intervalo \textbf{$r\in[h,h+w]$}:
\begin{align*}
    \Phi_B
    &=\int_h^{h+w}\frac{\mu_0 I}{2\pi r}L\,dr \\
    &=\frac{\mu_0 I L}{2\pi}\int_h^{h+w}\frac{dr}{r} \\
    &=\frac{\mu_0 I L}{2\pi}\left[\ln r\right]_h^{h+w} \\
    &=\frac{\mu_0 I L}{2\pi}\ln\left(\frac{h+w}{h}\right).
\end{align*}

Por tanto,
\begin{equation*}
    \boxed{\Phi_B=\frac{\mu_0 I L}{2\pi}\ln\left(\frac{h+w}{h}\right).}
\end{equation*}

\paso{Parte (b): fem inducida si $I=a+bt$.}
Como la corriente depende del tiempo,
\begin{equation*}
    I(t)=a+bt.
\end{equation*}

Sustituyendo en el flujo,
\begin{equation*}
    \Phi_B(t)=\frac{\mu_0 L}{2\pi}\ln\left(\frac{h+w}{h}\right)(a+bt).
\end{equation*}

Derivando respecto del tiempo,
\begin{equation*}
    \frac{d\Phi_B}{dt}=\frac{\mu_0 Lb}{2\pi}\ln\left(\frac{h+w}{h}\right).
\end{equation*}

Por Faraday,
\begin{equation*}
    \varepsilon=-\frac{d\Phi_B}{dt}.
\end{equation*}

Luego,
\begin{equation*}
    \boxed{\varepsilon=-\frac{\mu_0 Lb}{2\pi}\ln\left(\frac{h+w}{h}\right).}
\end{equation*}

Ahora pasamos los datos a unidades SI:
\begin{equation*}
    b=10\,\mathrm{A/s},\qquad
    h=1.0\times10^{-2}\,\mathrm{m},\qquad
    w=1.0\times10^{-1}\,\mathrm{m},\qquad
    L=1.0\,\mathrm{m}.
\end{equation*}

Adem\'as,
\begin{equation*}
    \frac{\mu_0}{2\pi}=2.0\times10^{-7}\,\mathrm{T\,m/A}.
\end{equation*}

Sustituyendo,
\begin{align*}
    \varepsilon
    &=-(2.0\times10^{-7})(1.0)(10)
    \ln\left(\frac{0.01+0.10}{0.01}\right)\,\mathrm{V} \\
    &=-2.0\times10^{-6}\ln(11)\,\mathrm{V}.
\end{align*}

Como $\ln(11)\approx2.40$,
\begin{equation*}
    \varepsilon\approx -4.8\times10^{-6}\,\mathrm{V}.
\end{equation*}

Entonces,
\begin{equation*}
    \boxed{\varepsilon\approx -4.8\,\mu\mathrm{V}.}
\end{equation*}

\paso{Direcci\'on de la corriente inducida.}
Si la corriente del alambre va hacia la derecha, la regla de la mano derecha indica que, en la regi\'on de la espira, el campo magn\'etico del alambre \textbf{entra al plano de la hoja}.

Como $b>0$, la corriente del alambre aumenta. Por tanto, el flujo entrando al plano tambi\'en aumenta:
\begin{equation*}
    \Phi_B^{\rm ext}\ \text{entrando al plano aumenta}.
\end{equation*}

Por la ley de Lenz, la espira debe crear un campo inducido saliendo del plano. Un campo saliendo del plano corresponde a una corriente antihoraria:
\begin{equation*}
    \boxed{I_{\rm ind}\ \text{circula en sentido antihorario}.}
\end{equation*}


Con los datos escritos en el enunciado, la fem queda del orden de microvoltios: $|\varepsilon|\approx4.8\,\mu\mathrm{V}$. Si una pauta entrega $4.8\,\mathrm{mV}$, probablemente hay una discrepancia en alguna unidad o dato num\'erico.


\end{document}
